\providecommand{\BM}{\textbf{[B]}}
\providecommand{\KM}{\textbf{[K]}}
\providecommand{\SM}{\textbf{[S]}}
\providecommand{\XM}{\textbf{[X]}}
\providecommand{\QM}{\textbf{[Q]}}

\chapter*{Schnittstellen der FFGFT zu anderen Rahmen}
\addcontentsline{toc}{chapter}{Schnittstellen der FFGFT zu anderen Rahmen}

{\small\noindent\textbf{Zusammenfassung.}\quad\ignorespaces
Der Korpus enthält Vergleichs- und Brückendokumente zu elf Rahmen, verteilt
über mehr als zwanzig Dokumente. Dieses Dokument fasst sie zusammen: je Rahmen
die Dokumentnummern, den mathematischen Berührungspunkt, den Status und den
offenen Punkt. Danach folgen die Größen, an denen sich mehrere Rahmen
tatsächlich überschneiden, jeweils mit Rechenweg, als Vorlage für eine
Vergleichstabelle. Es wird nichts Neues gerechnet; alle Zahlen stammen aus den
genannten Dokumenten und sind im Prüfskript nachgerechnet.
}\medskip

\section*{Zeichenerklärung}
\addcontentsline{toc}{section}{Zeichenerklärung}

\medskip
\begin{tabular}{@{}ll@{}}
\textbf{[B]} & algebraisch bewiesen \\
\textbf{[K]} & numerisch bestätigt \\
\textbf{[S]} & strukturell plausibel, Herleitung offen \\
\textbf{[X]} & geprüft und negativ \\
\textbf{[Q]} & aus externer Quelle übernommen, zitiert \\
\end{tabular}

\section{Grundsatz}

FFGFT leitet aus $\xi$ ab. Messwerte und Standardmodell-Größen dienen der
Umrechnung und der Prüfbarkeit, nicht der Ableitung; der kosmische Sektor
bleibt ausgeklammert (P39). Größen fremder Rahmen treten in allen unten
genannten Dokumenten als zitierte Eingänge \QM{} auf und stehen nicht
gleichrangig neben den eigenen Setzungen. Ein Berührungspunkt gilt erst dann
als Brücke, wenn dasselbe strukturelle Objekt aus der jeweils eigenen
Geometrie folgt und eine explizite Abbildung zwischen den Herleitungen
besteht; das Kriterium stammt aus Dok.~283.

\section{Übersicht}

\begin{center}
\small
\begin{tabular}{@{}llp{5.2cm}l@{}}
\toprule
\textbf{Rahmen} & \textbf{Dokumente} & \textbf{Berührungspunkt} & \textbf{Status} \\
\midrule
HLV (Krüger) & 271, 272, 276, 282, 283, 285, 294, 297 &
  Ikosaedrischer Cut-and-Project-Quasikristall gegen $T^4/\mathbb{Z}_3$;
  Dimensionsbrücke $6=\mathbf{3}\oplus\mathbf{3'}$ in $A_5$; CP-Teilbarkeit &
  Brücke \KM, Gabelung bei der Symmetrieordnung \\
OPH (Mueller) & 364, 374, 375, 376 &
  Spektrum des ikosaedrischen Trägers; Energie-Längen-Brücke
  $E_{\rm bit}=\hbar c/L$; $v/E_P$; Koeffizient $8\pi$ &
  Brückengleichungen \BM, Skalen offen auf Gegenseite \QM \\
Hyperbit, $G(6,\mathbb{Z}_3\mathbb{C})$ (Matzke) & 324, 326, 336 &
  $\mathbb{Z}_3$-Struktur und $\mathrm{GF}(9)$; Bit-Energie und
  Landauer-Schwelle &
  algebraische Brücke \BM, Bit-Energie systemabhängig \KM \\
XYLATIC (Phillips) & 252 &
  Zahlentheoretische Strukturen der $\sigma$-Funktion &
  Strukturvergleich \\
RA/PMT (Guevara) & 269 &
  Dreigliedriges Schichtenschema, Wahrnehmungsebene &
  Strukturvergleich \\
Vopson & 251 &
  Computational-Universe-Hypothese, zweites Gesetz der Infodynamik &
  verschiedene Primitive \\
Matsas et al. & 105 &
  Zahl der fundamentalen Konstanten &
  Einordnung \\
$\Lambda$CDM & 309 &
  Skalenanker: beide Rahmen brauchen genau eine Zahl zur SI-Kopplung &
  gemeinsame Lage \KM \\
Kagome-Gitter & 361 &
  Vier strukturelle Parallelen; Masse als geometrische Einfrierung &
  Analogie \SM \\
SM-Lagrangedichte & 049 &
  Feldzahl und freie Parameter gegen eine Gleichung &
  Formvergleich \\
Literatur allgemein & 345 &
  Einordnung gegenüber veröffentlichten Ansätzen &
  Einordnung \\
\bottomrule
\end{tabular}
\end{center}

UIFT (Teker) wird in Dok.~013 und 257 über die Bit-Masse berührt, hat aber
noch kein eigenes Vergleichsdokument.

\section{Größen, an denen sich Rahmen überschneiden}

Diese Tabelle ist die Vorlage für eine Vergleichstabelle. Zu jedem Wert gehört
der Rechenweg; ohne ihn ist der Wert nicht vergleichbar (Dok.~375).

\begin{center}
\small
\begin{tabular}{@{}llrl@{}}
\toprule
\textbf{Größe} & \textbf{FFGFT-Wert} & \textbf{Abw.} & \textbf{Weg} \\
\midrule
$\alpha^{-1}$ & $3700/27=137{,}037$ & $+7{,}6$\,ppm &
  Geometrie aus $T^4/\mathbb{Z}_3$ (Dok.~007), Galois-Bestätigung (Dok.~338) \\
$\sin^2\theta_W$ (on-shell) & $2/9=0{,}22222$ & $-0{,}52\,\%$ &
  $A_5$-Struktur (Dok.~293, 372) \\
$\lambda_{\rm CKM}$ & $\xi^{1/6}=0{,}22602$ & $+0{,}46\,\%$ &
  Galois-Brücke (Dok.~336) \\
$v/E_P$ & $2{,}0389\times10^{-17}$ & $+1{,}10\,\%$ &
  $v$ aus $\xi$ (R104), $E_P$ aus eigener Kette (Dok.~375) \\
$E_{\rm bit}(L)$ & $\hbar c/L$ & --- &
  systemabhängig, kein universeller Bitwert (Dok.~257, 326) \\
Dimension & $3{+}1$ aus $T^4/\mathbb{Z}_3$ & --- &
  Setzung; $6=\mathbf{3}\oplus\mathbf{3'}$ als Brücke zu HLV (Dok.~285) \\
Kosmischer Sektor & --- & --- &
  ausgeklammert (P39, Dok.~309) \\
\bottomrule
\end{tabular}
\end{center}

\section{Was daraus folgt}

\begin{itemize}\itemsep2pt
\item Zwei Rahmen sind über mehr als eine Größe verbunden: HLV über
  $A_5$-Struktur und Dimension, OPH über Träger, Skala und $v/E_P$. Die
  übrigen sind Einzelvergleiche.
\item Die Überschneidungen liegen fast alle im Leptonsektor und in der
  Geometrie. Der Hadronsektor kommt in keinem Vergleich vor; das ist keine
  Auswahl, sondern der Stand der Gegenseiten.
\item $\alpha^{-1}$ und $v/E_P$ sind derzeit die einzigen Größen, für die
  mehr als ein Rahmen eine Zahl angibt. Sie sind damit die einzigen Zeilen
  einer Vergleichstabelle, die Daten auseinanderhalten können.
\item Der Skalenanker ist bei jedem Rahmen dieselbe Lage an anderer Stelle
  (Dok.~309); er ist kein Unterscheidungsmerkmal.
\end{itemize}

\section*{Werkzeug und Methode}
\addcontentsline{toc}{section}{Werkzeug und Methode}
Claude (Anthropic) wurde als Werkzeug eingesetzt für: Zusammenstellen der
Dokumentliste, Schreiben und Ausführen des Prüfskripts, Ausarbeiten der
\LaTeX-Quellen und Pflege des Changelogs. Das Werkzeug rechnet und schreibt;
der Autor entscheidet, was berechnet und geschrieben wird, beurteilt jedes
Ergebnis und bestimmt, was es physikalisch bedeutet. Das Prüfskript
kontrolliert, dass jedes zitierte Dokument im Korpus vorliegt, und rechnet
jede Zahl der zweiten Tabelle aus FFGFT-Größen nach.

\section{Ergebnis}

\begin{enumerate}
\item Elf Rahmen mit Vergleichs- oder Brückendokumenten, davon zwei mit
  mehreren Berührungspunkten \KM.
\item Sieben Größen, an denen sich Rahmen überschneiden, jeweils mit Weg und
  Abweichung \KM.
\item Nur $\alpha^{-1}$ und $v/E_P$ werden von mehr als einem Rahmen
  beziffert.
\item Der Hadronsektor kommt in keinem Vergleich vor.
\end{enumerate}

\medskip
\noindent\textbf{Prüfskript:} \texttt{2/python/Dok377\_Skripte/pruef\_377\_schnittstellen.py} (20/20)

\medskip
\noindent\textbf{Register (Dok.~190):} kein Eintrag.

\begin{thebibliography}{9}
\bibitem{hlv} Dok.~271, 272, 276, 282, 283, 285, 294, 297 (HLV).
\bibitem{oph} Dok.~364, 374, 375, 376 (OPH).
\bibitem{hyperbit} Dok.~324, 326, 336 (Hyperbit, $\mathbb{Z}_3\mathbb{C}$).
\bibitem{weitere} Dok.~049, 105, 251, 252, 269, 309, 345, 361.
\bibitem{quellen} Dok.~007, 013, 257, 293, 338, 372 (FFGFT-Quellen der Zahlen).
\end{thebibliography}
